\chapter{Parseur syntaxique}
\label{chap:parseur}

\section{Langage de requêtes WIDO}

Le langage \wido{} permet d'exprimer des requêtes formelles à variables sur le graphe JDM. Sa grammaire (simplifiée) est la suivante :

\begin{verbatim}
requête  ::= expr
expr     ::= terme (ET terme)*
terme    ::= facteur (OU facteur)*
facteur  ::= '(' clause ')'
           | '(' expr ')'
clause   ::= '(' sujet relation objet ')'
           | '(' variable '=' filtre ')'
sujet    ::= variable | constante
objet    ::= variable | constante
variable ::= '$' identifiant
filtre   ::= [a-z%]+
\end{verbatim}

\section{Exemples de requêtes}

\begin{lstlisting}[style=widoquery, caption=Exemples de requêtes WIDO]
($x r_isa animal)
($x r_isa animal) ET ($x = ch%)
($x r_isa artiste) ET (($x = ba%) OU ($x = Ba%))
($x r_isa animal) ET ($y r_isa animal) ET ($x r_can_eat $y)
(chat r_has_part $y) ET ($y r_isa $z)
\end{lstlisting}

\section{Structure de l'AST}

Le parseur produit un arbre de décision récursif. Chaque nœud interne représente un opérateur logique (ET/OU), chaque feuille est une clause atomique (relation ou filtre).

\begin{lstlisting}[caption=AST produit pour \texttt{(\$x r\_isa animal) ET (\$x = ch\%)}]
{
  "type": "NOEUD_LOGIQUE",
  "operateur": "ET",
  "gauche": {
    "type": "CLAUSE_RELATION",
    "variable": "$x",
    "relation": "r_isa",
    "cible": "animal"
  },
  "droite": {
    "type": "CLAUSE_FILTRE",
    "variable": "$x",
    "filtre": "ch%"
  }
}
\end{lstlisting}

\section{Gestion des erreurs syntaxiques}

Lorsqu'une clause est incomplète (ex: \texttt{(\$x r\_isa)}), le parseur lève une exception qui est capturée par le serveur. Celui-ci retourne une réponse JSON propre avec \code{statut: "Erreur"} et un message explicatif, sans jamais retourner de HTTP 500 ni de champs \code{undefined}.
